%% Compiler: XeLaTeX (Overleaf: Menu -> Compiler -> XeLaTeX)
\documentclass[12pt]{article}
\usepackage[UTF8]{ctex}
\usepackage[margin=1in]{geometry}
\usepackage{booktabs,array,longtable,hyperref,enumitem,xurl}
\onehalfspacing
\newcommand{\path}[1]{\texttt{\detokenize{#1}}}
\title{Machine Learning and Financial Technology --- Final Project}
\author{112700019 Paul}
\date{Fall 2026}
\begin{document}
\maketitle
\section{Paper Replication}
Feng, Fuli, Xiangnan He, Xiang Wang, Cheng Luo, Yiqun Liu, and Tat-Seng Chua.
``Temporal Relational Ranking for Stock Prediction.'' \textit{ACM Transactions on Information Systems} 37, no.~2 (2019): Article 27.
\url{https://doi.org/10.1145/3309547}.
\begin{itemize}[leftmargin=*]
\item Journal tier: entry 11 in the NYCU CS A-tier journal list (\path{journal-ranking/交大資工A級期刊_20200910Updated.xlsx}).
\item Official code: \url{https://github.com/fulifeng/Temporal_Relational_Stock_Ranking}.
\item Set \path{TF_USE_LEGACY_KERAS=1} before training. Keras 3 does not provide \path{BasicLSTMCell}; see \path{reports/R3_model_and_experiments/PROGRESS.md}.
\end{itemize}
\section{Paper Summary and Method}
The paper proposes Relational Stock Ranking (RSR) with Temporal Graph Convolution, encoding industry and Wikidata relations among stocks into a time-sensitive ranking task to predict the relative ranking of next-day stock returns.
\section{What Is Actually Graded}
The R1--R4 folders originally referred to a grading-rubric item for which no source could be found. Neither course syllabus nor the official course GitHub contains that rubric or describes separate grades for R1, R2, R3, and R4.
\begin{center}\begin{tabular}{lll}\toprule
Item & Weight & Description \\\midrule
Participation & 20\% & Exercises, presentations, class wrap-ups, and cold calls \\
Project & 40\% & LaTeX manuscript and presentation slides \\
Exam & 40\% & In-class exam; one double-sided A4 reference sheet allowed \\\bottomrule
\end{tabular}\end{center}
The Project grade is based on the LaTeX manuscript, not the folder structure. R1--R4 hold draft material on motivation, EDA, methods, experiments, results, and conclusions.
\section{Project Scope}
This project is a straight replication of Feng et al. (2019), without an architectural extension. It compares Rank\_LSTM (without a relation graph) with RSR (using a static industry graph). Dynamic-relation experiments are retained as exploratory records and excluded from the final manuscript's conclusions.
\section{Data and Code}
Processed stock data beginning 2013-01-01 are under \path{../Temporal_Relational_Stock_Ranking/data/}. Industry relations encode sector/industry links; Wikidata relations encode company relationships.
\begin{itemize}[leftmargin=*]
\item \path{training/rank_lstm.py}: Rank\_LSTM baseline without a relation graph.
\item \path{training/relation_rank_lstm.py}: Relational Stock Ranking (RSR).
\end{itemize}
\section{Manuscript and Slides}
\begin{itemize}[leftmargin=*]
\item Overleaf: \url{https://www.overleaf.com/read/fjkrqhnbqwgc#e5d02b} (read-only link; add \path{venteng@gmail.com} as a collaborator).
\item Canva slides: \url{https://canva.link/psyac0w98zuyn5l}.
\item GitHub: \url{https://github.com/paul931130/112700019-Paul}.
\end{itemize}
\end{document}